\begin{abstract}
在 API 计费与上下文长度约束下，函数级代码生成的部署效果取决于测试时计算如何组织，而非单次生成本身。本文将问题表述为可复现的\textbf{预算分档编排协议}：在单一实现与统一统计口径下，把从单步到级联的行为固定为四个可对照档位，使不同预算形态可在同一执行器与同一评测口径下比较。

主实验在 HumanEval 与 MBPP 两条函数级基准上联合报告通过率、调用次数与总 token，并给出经验收益率 $\rho$ 刻画各档边际与累计成本--收益关系。结果表明：在更难、更杂的任务分布上，分档编排带来更宽的通过率提升空间；在易饱和分布上，增益更多体现为残差段收束与成本重分配。机制上，静态结构难度分更适合作为\textbf{预算分配先验}而非运行失败分类器；multiset 随机路由对照用于分离 hard 覆盖率与结构对齐的贡献；反馈回注形态对照表明，压缩回注在精度几乎不变的前提下可显著降低 token 开销。

部署上，中预算侧重结构化主链与失败回注修补，高预算侧重候选扩展并在可选 refine 上做精度与 token 的折中。
\keywords{函数级代码生成；测试时计算；预算分档；失败修补；级联编排}
\end{abstract}

\begingroup
\renewcommand{\abstractname}{Abstract}
\renewcommand{\abstracttextfont}{\small\normalfont\noindent\ignorespaces}
\begin{abstract}
Under API pricing and context-length constraints, the deployment effectiveness of function-level code generation depends on how test-time computation is organized, rather than on a single generation attempt alone. This paper formulates the problem as a reproducible \textbf{budget-tiered orchestration protocol}: under one implementation and one unified accounting interface, behaviors from single-step generation to cascaded workflows are fixed into four comparable tiers, so that different budget forms can be compared under the same executor and evaluation protocol.

The main experiments jointly report pass rate, number of calls, and total tokens on the HumanEval and MBPP function-level benchmarks, and introduce an empirical gain rate $\rho$ to characterize both marginal and cumulative cost--benefit relationships across tiers. Results show that tiered orchestration provides broader room for pass-rate improvement on harder and more heterogeneous task distributions; on easier and more saturated distributions, the gains are more reflected in residual-set convergence and cost reallocation. Mechanistically, static structural difficulty scores are more suitable as \textbf{budget-allocation priors} than as runtime failure classifiers; the multiset random-routing control separates the contribution of hard-slot coverage from structural alignment; and the feedback-injection comparison shows that compressed feedback can substantially reduce token cost while keeping accuracy nearly unchanged.

For deployment, the medium-budget setting emphasizes a structured main chain and failure-feedback repair, whereas the high-budget setting emphasizes candidate expansion and trades off accuracy against token cost through optional refine steps.

\par\vskip2ex\noindent\normalfont{\bfseries Keywords: } function-level code generation; test-time computation; budget tiering; failure repair; cascade orchestration
\end{abstract}
\endgroup
